% GENERATED FILE. Edit canonical lessons, then run npm run generate:books.
% canonical-source-hash: fnv1a64:dbb38b7712f2db47
% canonical-lessons: TA-C414-vica, TA-C414-ellai, TA-C414-cavati, TA-C414-nulaivayil, TA-C414-parimutal-cey

\chapter{At the Border}
\label{ch:ta-at-the-border}
\hlchaptermodality{414}

\begin{chapteropening}
\textbf{By the end of this chapter:} \emph{I can talk about a visa, a border, a checkpoint and an entrance, and goods being seized.}

Complete the last lesson of chapter 414: I can talk about a visa, a border, a checkpoint and an entrance, and goods being seized.
\end{chapteropening}

\section[\texorpdfstring{\ta{விசா} (vicā)}{விசா (vicā)}]{\ta{விசா} (vicā) — a visa}

\label{lesson:TA-C414-vica}



\begin{quote}
Before the new one: say the Tamil for customs, then the Tamil for immigration.
\end{quote}

\subsection*{You'll want to know: \ta{விசா}}
\textbf{\ta{விசா}} — \emph{vicā} — \textquotedblleft{}a visa\textquotedblright{}.

\begin{quote}
\ta{அந்த} \ta{நாட்டுக்குப்} \ta{போக} \ta{விசா} \ta{வேண்டும்}.
\end{quote}

\begin{quote}
\emph{anta nāṭṭukkup pōka vicā vēṇṭum.}
\end{quote}

\begin{quote}
You need a visa to go to that country.
\end{quote}

\ta{விசா} \ta{வாங்கு} is to get a visa.

\begin{cousinweb}
\textbf{\ta{விசா}} (\emph{vicā}) is a loan from English \emph{visa}.
\end{cousinweb}

\begin{grammarlens}[title={What you've built}]
One more word for getting around.
\end{grammarlens}

\subsection*{Your turn}
\begin{itemize}
\raggedright
  \item \emph{Say it:} \emph{vicā}
  \item \emph{Say it:} \emph{vicā}, once more
  \item \emph{Say it:} \emph{anta nāṭṭukkup pōka vicā vēṇṭum.} — You need a visa to go to that country.
  \item \emph{Recall:} say the Tamil for customs, then the Tamil for immigration, then say \emph{vicā} again
\end{itemize}

\begin{tcolorbox}[breakable,colback=teal!4,colframe=teal!35!black,title={Before you move on}]
What does \ta{விசா} mean? (\textquotedblleft{}A visa\textquotedblright{}.) Say it once more.
\end{tcolorbox}
\section[\texorpdfstring{\ta{எல்லை} (ellai)}{எல்லை (ellai)}]{\ta{எல்லை} (ellai) — a border; a limit}

\label{lesson:TA-C414-ellai}



\begin{quote}
Before the new one: say the Tamil for immigration, then the Tamil for a visa.
\end{quote}

\subsection*{You'll want to know: \ta{எல்லை}}
\textbf{\ta{எல்லை}} — \emph{ellai} — \textquotedblleft{}a border; a limit\textquotedblright{}.

\begin{quote}
\ta{எல்லையில்} \ta{நீண்ட} \ta{வரிசை} \ta{இருந்தது}.
\end{quote}

\begin{quote}
\emph{ellaiyil nīṇṭa varicai iruntatu.}
\end{quote}

\begin{quote}
There was a long queue at the border.
\end{quote}

\ta{எல்லைக்} \ta{கோடு} is a boundary line.

\begin{cousinweb}
\textbf{\ta{எல்லை}} (\emph{ellai}) is native Dravidian (DEDR 846), \textquotedblleft{}a limit, a border, a boundary\textquotedblright{}.
\end{cousinweb}

\begin{grammarlens}[title={What you've built}]
One more word for getting around.
\end{grammarlens}

\subsection*{Your turn}
\begin{itemize}
\raggedright
  \item \emph{Say it:} \emph{ellai}
  \item \emph{Say it:} \emph{ellai}, once more
  \item \emph{Say it:} \emph{ellaiyil nīṇṭa varicai iruntatu.} — There was a long queue at the border.
  \item \emph{Recall:} say the Tamil for immigration, then the Tamil for a visa, then say \emph{ellai} again
\end{itemize}

\begin{tcolorbox}[breakable,colback=teal!4,colframe=teal!35!black,title={Before you move on}]
What does \ta{எல்லை} mean? (\textquotedblleft{}A border; a limit\textquotedblright{}.) Say it once more.
\end{tcolorbox}
\section[\texorpdfstring{\ta{சாவடி} (cāvaṭi)}{சாவடி (cāvaṭi)}]{\ta{சாவடி} (cāvaṭi) — a checkpoint booth; a booth}

\label{lesson:TA-C414-cavati}



\begin{quote}
Before the new one: say the Tamil for a visa, then the Tamil for a border.
\end{quote}

\subsection*{You'll want to know: \ta{சாவடி}}
\textbf{\ta{சாவடி}} — \emph{cāvaṭi} — \textquotedblleft{}a checkpoint booth; a booth\textquotedblright{}.

\begin{quote}
\ta{சுங்கச்} \ta{சாவடியில்} \ta{வண்டிகள்} \ta{நின்றன}.
\end{quote}

\begin{quote}
\emph{cuṅkac cāvaṭiyil vaṇṭikaḷ niṉṟaṉa.}
\end{quote}

\begin{quote}
The vehicles stopped at the toll booth.
\end{quote}

\ta{சுங்கச்} \ta{சாவடி} is a toll booth.

\begin{cousinweb}
\textbf{\ta{சாவடி}} (\emph{cāvaṭi}), \textquotedblleft{}a booth, a rest house\textquotedblright{}, is an old Tamil word whose origin is not settled.
\end{cousinweb}

\begin{grammarlens}[title={What you've built}]
One more word for getting around.
\end{grammarlens}

\subsection*{Your turn}
\begin{itemize}
\raggedright
  \item \emph{Say it:} \emph{cāvaṭi}
  \item \emph{Say it:} \emph{cāvaṭi}, once more
  \item \emph{Say it:} \emph{cuṅkac cāvaṭiyil vaṇṭikaḷ niṉṟaṉa.} — The vehicles stopped at the toll booth.
  \item \emph{Recall:} say the Tamil for a visa, then the Tamil for a border, then say \emph{cāvaṭi} again
\end{itemize}

\begin{tcolorbox}[breakable,colback=teal!4,colframe=teal!35!black,title={Before you move on}]
What does \ta{சாவடி} mean? (\textquotedblleft{}A checkpoint booth; a booth\textquotedblright{}.) Say it once more.
\end{tcolorbox}
\section[\texorpdfstring{\ta{நுழைவாயில்} (nuḻaivāyil)}{நுழைவாயில் (nuḻaivāyil)}]{\ta{நுழைவாயில்} (nuḻaivāyil) — an entrance; a gateway}

\label{lesson:TA-C414-nulaivayil}



\begin{quote}
Before the new one: say the Tamil for a border, then the Tamil for a checkpoint booth.
\end{quote}

\subsection*{You'll want to know: \ta{நுழைவாயில்}}
\textbf{\ta{நுழைவாயில்}} — \emph{nuḻaivāyil} — \textquotedblleft{}an entrance; a gateway\textquotedblright{}.

\begin{quote}
\ta{நுழைவாயிலில்} \ta{காவலர்கள்} \ta{நின்றார்கள்}.
\end{quote}

\begin{quote}
\emph{nuḻaivāyilil kāvalarkaḷ niṉṟārkaḷ.}
\end{quote}

\begin{quote}
Guards stood at the entrance.
\end{quote}

\ta{நுழை}, to enter, + \ta{வாயில்}, a doorway.

\begin{cousinweb}
\textbf{\ta{நுழைவாயில்}} (\emph{nuḻaivāyil}) is built from \ta{நுழை}, \textquotedblleft{}to enter\textquotedblright{} (DEDR 3714), and \ta{வாயில்}, \textquotedblleft{}a doorway, a gate\textquotedblright{} (DEDR 5354).
\end{cousinweb}

\begin{grammarlens}[title={What you've built}]
One more word for getting around.
\end{grammarlens}

\subsection*{Your turn}
\begin{itemize}
\raggedright
  \item \emph{Say it:} \emph{nuḻaivāyil}
  \item \emph{Say it:} \emph{nuḻaivāyil}, once more
  \item \emph{Say it:} \emph{nuḻaivāyilil kāvalarkaḷ niṉṟārkaḷ.} — Guards stood at the entrance.
  \item \emph{Recall:} say the Tamil for a border, then the Tamil for a checkpoint booth, then say \emph{nuḻaivāyil} again
\end{itemize}

\begin{tcolorbox}[breakable,colback=teal!4,colframe=teal!35!black,title={Before you move on}]
What does \ta{நுழைவாயில்} mean? (\textquotedblleft{}An entrance; a gateway\textquotedblright{}.) Say it once more.
\end{tcolorbox}
\section[\texorpdfstring{\ta{பறிமுதல்} \ta{செய்} (paṟimutal cey)}{பறிமுதல் செய் (paṟimutal cey)}]{\ta{பறிமுதல்} \ta{செய்} (paṟimutal cey) — to confiscate; to seize}

\label{lesson:TA-C414-parimutal-cey}



\begin{quote}
Before the new one: say the Tamil for a checkpoint booth, then the Tamil for an entrance.
\end{quote}

\subsection*{You'll want to know: \ta{பறிமுதல்} \ta{செய்}}
\textbf{\ta{பறிமுதல்} \ta{செய்}} — \emph{paṟimutal cey} — \textquotedblleft{}to confiscate; to seize\textquotedblright{}.

\begin{quote}
\ta{அதிகாரிகள்} \ta{அவனுடைய} \ta{பையைப்} \ta{பறிமுதல்} \ta{செய்தார்கள்}.
\end{quote}

\begin{quote}
\emph{atikārikaḷ avaṉuṭaiya paiyaip paṟimutal ceytārkaḷ.}
\end{quote}

\begin{quote}
The officials confiscated his bag.
\end{quote}

Past (I …): \ta{பறிமுதல்} \ta{செய்தேன்}, \emph{paṟimutal ceytēṉ}.

What is seized takes -\ta{ஐ}: \ta{பையைப்} \ta{பறிமுதல்} \ta{செய்}.

\begin{cousinweb}
\textbf{\ta{பறிமுதல்} \ta{செய்}} (\emph{paṟimutal cey}) joins the noun \ta{பறிமுதல்}, \textquotedblleft{}confiscation\textquotedblright{}, whose origin is not settled, to \ta{செய்}, \textquotedblleft{}to do\textquotedblright{}.
\end{cousinweb}

\begin{grammarlens}[title={What you've built}]
One more word for getting around.
\end{grammarlens}

\subsection*{Your turn}
\begin{itemize}
\raggedright
  \item \emph{Say it:} \emph{paṟimutal cey}
  \item \emph{Say it:} \emph{paṟimutal cey}, once more
  \item \emph{Say it:} \emph{atikārikaḷ avaṉuṭaiya paiyaip paṟimutal ceytārkaḷ.} — The officials confiscated his bag.
  \item \emph{Recall:} say the Tamil for a checkpoint booth, then the Tamil for an entrance, then say \emph{paṟimutal cey} again
\end{itemize}

\begin{tcolorbox}[breakable,colback=teal!4,colframe=teal!35!black,title={Before you move on}]
What does \ta{பறிமுதல்} \ta{செய்} mean? (\textquotedblleft{}To confiscate; to seize\textquotedblright{}.) Say it once more.
\end{tcolorbox}
